\documentclass[UTF8,zihao=-4]{ctexart}
\usepackage[a4paper,margin=2.2cm]{geometry}
\usepackage{xcolor}
\usepackage{hyperref}
\usepackage{enumitem}
\usepackage{fontspec}
\usepackage{amsmath}
\IfFontExistsTF{Microsoft YaHei}{\setCJKmainfont{Microsoft YaHei}}{\setCJKmainfont{SimSun}[BoldFont=SimHei]}
\IfFontExistsTF{Microsoft YaHei UI}{\setCJKsansfont{Microsoft YaHei UI}}{}
\IfFontExistsTF{Consolas}{\setmonofont{Consolas}}{}
\definecolor{linkblue}{RGB}{30,80,145}
\hypersetup{colorlinks=true,linkcolor=linkblue,urlcolor=linkblue,pdfborder={0 0 0}}
\setlength{\parindent}{2em}
\setlength{\parskip}{0.25em}
\linespread{1.16}
\setlist{itemsep=0.2em,topsep=0.35em,leftmargin=2.5em}
\pagestyle{plain}

\title{07\_对抗策略分析}

\author{}

\date{}

\begin{document}

\maketitle

\section*{07 对抗策略分析}

\subsection*{题目原文}

能解释对抗策略的含义、能采用对抗策略来分析经典问题：同时求最大和最小的问题和求取第二大元素的问题。

\subsection*{考查类型}

概念题和证明题。

\subsection*{对抗策略含义}

对抗策略用于证明问题的下界。分析者假设存在一个对手，它在算法每次询问或比较时给出回答，并保持这些回答与至少一个合法输入一致。

如果对手能让任何算法都必须进行至少 $L(n)$ 次操作才能确定答案，则得到下界：

\[\Omega(L(n))\]

这种证明不是设计一个慢算法，而是证明所有算法都避不开一定数量的比较。

\subsection*{同时求最大值和最小值}

朴素做法分别找最大和最小，需要：

\[(n - 1) + (n - 1) = 2n - 2\]

更优做法是成对比较。

步骤：

\begin{enumerate}
\item 每两个元素先相互比较一次。
\item 每对中较大的元素进入最大值比较范围。
\item 每对中较小的元素进入最小值比较范围。
\item 分别在两个范围内找最大和最小。
\end{enumerate}

当 $n$ 为偶数时，比较次数为：

\[n/2 + (n/2 - 1) + (n/2 - 1) = 3n/2 - 2\]

当 $n$ 为奇数时，可以先留出一个元素作为初始最大值和最小值，再对剩余 \texttt{n - 1} 个元素成对处理。比较次数为：

\begin{quote}
\small\ttfamily\raggedright
\relax 3(n\ -\ 1)/2\\
\end{quote}

统一写法：

\[\lceil 3n/2 \rceil - 2\]

\subsection*{下界证明思路}

每个元素在比较历史中需要获得足够信息：

一个元素只有输过一次，才可以排除它成为最大值。

一个元素只有赢过一次，才可以排除它成为最小值。

第一次参与比较时，一次比较能同时给两个元素分配角色：胜者不再是最小值，败者不再是最大值。之后要确定最大值和最小值，两个集合内还要分别完成淘汰。

由此得到同时求最大值和最小值的比较下界，与成对算法的次数匹配。

\subsection*{求第二大元素}

先用锦标赛方式找最大值。每次比较的胜者继续向上，最终最大值会胜出。

找最大值需要：

\begin{quote}
\small\ttfamily\raggedright
\relax n\ -\ 1\\
\end{quote}

次比较。

第二大元素只会出现在直接输给最大值的元素中。若最大值直接比较并击败了 \texttt{k} 个元素，则还需要在这 \texttt{k} 个元素中找最大值，比较次数为：

\begin{quote}
\small\ttfamily\raggedright
\relax k\ -\ 1\\
\end{quote}

通过平衡锦标赛组织比较，可以使：

\[k = \lceil \log_2 n \rceil\]

总比较次数为：

\[n - 1 + \lceil \log_2 n \rceil - 1\]
\[= n + \lceil \log_2 n \rceil - 2\]

\subsection*{第二大元素的下界思路}

要确认最大值，至少需要 \texttt{n - 1} 次比较，因为除了最大值外的每个元素都必须至少输一次。

要确认第二大元素，只需在直接输给最大值的元素中查找。对抗策略可使最大值必须直接击败至少 $\lceil \log_2 n \rceil$ 个元素，因此还需要至少 $\lceil \log_2 n \rceil - 1$ 次比较。

所以比较下界为：

\[n + \lceil \log_2 n \rceil - 2\]

\subsection*{例题}

\subsubsection*{例题 1}

题目：$n = 10$ 时，同时求最大值和最小值最少需要多少次比较？给出达到该次数的策略。

解答：

比较次数为：

\[\lceil 3n/2 \rceil - 2 = \lceil 15 \rceil - 2 = 13\]

策略：

\begin{enumerate}
\item 将 10 个元素两两配对，先比较 5 次。
\item 每对较大的 5 个元素中求最大值，需要 4 次比较。
\item 每对较小的 5 个元素中求最小值，需要 4 次比较。
\end{enumerate}

总次数：

\[5 + 4 + 4 = 13\]

\subsubsection*{例题 2}

题目：用锦标赛方法在下面数组中求第二大元素，并写出比较次数。

\begin{quote}
\small\ttfamily\raggedright
\relax [9,\ 4,\ 7,\ 12,\ 3,\ 10,\ 6,\ 8]\\
\end{quote}

解答：

第一轮：

\[9 > 4\]
\[12 > 7\]
\[10 > 3\]
\[8 > 6\]

第二轮：

\[12 > 9\]
\[10 > 8\]

第三轮：

\[12 > 10\]

最大值是 \texttt{12}。直接输给 \texttt{12} 的元素是：

\begin{quote}
\small\ttfamily\raggedright
\relax 7,\ 9,\ 10\\
\end{quote}

在这三个元素中再找最大值，需要 2 次比较，得到第二大元素：

\begin{quote}
\small\ttfamily\raggedright
\relax 10\\
\end{quote}

总比较次数：

\[7 + 2 = 9\]

公式验证：

\[n + \lceil \log_2 n \rceil - 2 = 8 + 3 - 2 = 9\]

\subsection*{常见误区}

第二大元素不是在所有非最大元素中重新找最大。只需要检查直接输给最大值的元素。

对抗策略证明的是任何算法都需要的最低比较次数，不是某个具体算法的运行过程。

\end{document}
